\documentclass[../../../thesis.tex]{subfiles}
\begin{document}
\subsection{Webová stránka, sociálna sieť, video}
Pri multimediálnom online obsahu často nepoznáme autora,
prípadne je komplikované zistiť názov webovej stránky.
Snažíme sa teda zahrnúť čo najviac jednoznačných informácií,
najmä webovú adresu, dátum publikovania a dátum citovania.
\begin{trivlist}
  \item TVORCOVIA. \textit{Názov obsahu} [online]. Platforma, dátum publikovania [cit. dátum citovania]. Dostupnosť.
\end{trivlist}

\subsubsection*{\normalsize Príklady}
\begin{enumerate}
  \item VALKO, P. \textit{M31, M32 a jeden Starlink k tomu} [online]. Facebook, 2024-08-30 [cit. 2024-09-06]. Dostupné z: \texttt{https://www.facebook.com/share/p/S93xZoz9Rbnh
  dQ7M}.

  \item \textit{Why lenses can't make perfect images} [online]. Youtube, 2017-10-18 [cit. 2024-08-25]. Dostupné z: \texttt{https://youtu.be/DDoryfCXxPI?si=hC5Kuuf4p3OxOOsm}.
\end{enumerate}

\subsubsection*{\normalsize Zápis v zdrojovom .bib súbore}
\begin{verbatim}
@online{Valko2024M31,
  author       = {Pavol Valko},
  title        = {31, M32 a jeden Starlink k tomu},
  howpublished = {online},
  date         = {2025-08-30T00:03:00},
  url          = {https://www.facebook.com/share/p/S93xZoz9RbnhdQ7M},
  urldate      = {2024-09-06},
  publisher    = {Facebook},
}

@online{WhyLenses2017,
  title        = {Why lenses can’t make perfect images},
  howpublished = {online},
  date         = {2025-10-18},
  url          = {https://youtu.be/DDoryfCXxPI?si=hC5Kuuf4p3OxOOsm},
  urldate      = {2024-08-25},
  publisher    = {Youtube},
}
\end{verbatim}
\end{document}